\chapter{INTRODUCTION}

The electrocardiogram (ECG) is a non-invasive recording of the electrical activity of the heart, measured as potential differences between electrodes placed on the limbs and the chest. The standard 12-lead ECG combines limb leads and precordial leads to observe cardiac electrical activity from several directions, and it remains the reference recording for assessing rhythm, conduction and myocardial injury \cite{KLIGFIELD2007AHA}. Many diagnostic findings are tied to particular leads, and the AHA/ACCF/HRS recommendations set out separate criteria for intraventricular conduction disturbances, such as bundle branch block, and for acute ischaemia and infarction \cite{SURAWICZ2009AHA,WAGNER2009AHA}. The diagnostic value of an ECG therefore depends on which leads were recorded.

Automated interpretation of the ECG has developed from rule-based measurement programs \cite{SCHLAPFER2017COMPUTER} into deep neural networks trained on large clinical databases \cite{HANNUN2019CARDIOLOGIST,STRODTHOFF2021PTBXL}. More recently, ECG foundation models have been pretrained on very large collections of unlabelled recordings using self-supervised objectives, so that a single pretrained encoder can provide general-purpose signal representations for many downstream tasks \cite{MCKEEN2025ECGFM}. In practice, however, a complete 12-lead recording is not always available. Smaller and easier-to-use recording devices often capture only a subset of the leads, and single-lead ambulatory monitors are already in clinical use \cite{HANNUN2019CARDIOLOGIST,REYNA2022ISSUES}. When leads are missing, a model may still produce a confident diagnosis, but the physiological evidence that a clinician would require for that diagnosis may no longer be present in the signal.

This project, titled Reasoning Failure Boundaries in ECG Foundation Models Under Progressive Physiological Information Loss, studies this situation in a controlled manner. ECG recordings from three public cohorts, PTB-XL, Chapman-Shaoxing and MIMIC-IV-ECG, are standardised to a common format and evaluated on eight conditions: normal ECG, atrial fibrillation, right and left bundle branch block, sinus bradycardia, sinus tachycardia, first-degree atrioventricular block and myocardial infarction. Each recording is presented at seven lead-availability levels of 12, 10, 8, 6, 4, 2 and 1 leads. Missing leads are set to zero and are never reconstructed, so that the model only receives information that was actually recorded. A publicly released ECG foundation model, ECG-FM, is used as a frozen encoder that converts each degraded recording into a 768-dimensional representation \cite{MCKEEN2025ECGFM}.

On top of this encoder, the project builds a reasoning stage that combines the learned representation with rule-based measurements extracted only from the available leads. A Control Reasoning Transformer, which uses no external clinical knowledge, has been implemented and trained during the first half of the project. It produces probabilities for the eight conditions together with an attention-based reasoning trace. The project further proposes a deterministic grounding gate that labels each clinical claim as supported, contradicted or unavailable, and a knowledge-augmented reasoning path based on a Clinical Knowledge Prior derived from SNOMED CT. Both of these components have been specified but not yet implemented, and they are treated in this report as planned work.

\section{MOTIVATION}
The ECG is a quick and non-invasive cardiac investigation that is informative across a broad range of rhythm, conduction and ischaemic conditions \cite{KLIGFIELD2007AHA}. Its interpretation relies on measurements derived from the waveform, which lead to standardised diagnostic statements \cite{KLIGFIELD2007AHA}. Because these criteria refer to specific leads and measurements, an interpretation can in principle be checked against the recording from which it was made.

Automated ECG interpretation has a long history, and a review of computer-interpreted ECGs notes that the majority of the millions of ECGs recorded each year are analysed automatically, that limitations in diagnostic accuracy have persisted, and that clinicians may accept computer statements without verifying them \cite{SCHLAPFER2017COMPUTER}. Deep learning has since improved classification performance substantially on large datasets \cite{HANNUN2019CARDIOLOGIST,STRODTHOFF2021PTBXL}. As performance improves, it becomes more important to understand what evidence an automated system uses before its output is relied upon.

Foundation models extend this development by learning representations from unlabelled data that can be reused across tasks \cite{MCKEEN2025ECGFM}. ECG-FM, the model used in this project, follows the wav2vec 2.0 architecture and was pretrained with a combination of masked prediction, contrastive multi-segment coding and random lead masking \cite{MCKEEN2025ECGFM,OH2022LEADAGNOSTIC}. Random lead masking is intended to make a model usable when only some leads are present, but this raises a question that accuracy alone cannot answer: whether a diagnosis produced after leads are removed is still justified by the remaining signal.

Reduced-lead recordings are common in practice. Smaller and lower-cost devices that record fewer leads may widen access to cardiac care, although their diagnostic potential is not yet clear \cite{REYNA2022ISSUES}, and single-lead ambulatory monitors are already used for rhythm assessment \cite{HANNUN2019CARDIOLOGIST}. Research on reduced-lead classification has shown that for some tasks a limited set of leads can retain much of the classification performance of the full 12-lead ECG \cite{REYNA2022ISSUES,LAI2021LEADSELECTION}. Preserved classification performance, however, does not imply that the specific evidence for each diagnosis is still observable.

Progressive lead removal provides a controlled way to study this question. Because every recording is evaluated at each of the seven availability levels, any change in the model's behaviour can be attributed to the removed leads rather than to differences between patients. Missing leads are not reconstructed from the remaining ones. Reconstructed leads would be estimates rather than measurements \cite{PRESACAN2025RECON}. Masking without reconstruction keeps a clear distinction between information that was recorded and information that was not.

A further motivation concerns explanations. Current explanation methods are unlikely to provide reliable explanations for individual patients \cite{GHASSEMI2021FALSEHOPE}, and attention weights in particular cannot be assumed to explain a model's decision \cite{JAIN2019ATTENTION}. A useful explanation of an ECG diagnosis should therefore be checked against measurable evidence, and the check must separate three situations: the evidence is present and agrees with the claim, the evidence is present and disagrees with it, or the evidence cannot be assessed because the required leads are absent. This need to identify where model claims stop being supported by the available leads is the motivation for investigating Reasoning Failure Boundaries in ECG foundation models.

\section{BACKGROUND}
The standard 12-lead ECG is recorded from ten electrodes. Electrodes on the right arm, left arm and left leg form the limb leads I, II and III, the right-leg electrode serves as a reference, and six chest electrodes form the precordial leads V1 to V6. The augmented limb leads aVR, aVL and aVF are computed from the same limb electrodes \cite{KLIGFIELD2007AHA}. Consequently, lead III and the three augmented leads can be expressed as linear combinations of leads I and II, whereas the six precordial leads carry information that cannot be obtained from the limb leads. This relationship is important for the present work, since removing a derived limb lead does not remove electrical information that is available from leads I and II, whereas removing precordial leads does.

Different leads observe the heart from different directions, and diagnostic criteria are defined with respect to these views. The limb leads describe activity in the frontal plane, while the precordial leads describe activity in the horizontal plane. Separate parts of the AHA/ACCF/HRS recommendations address intraventricular conduction disturbances, including bundle branch block \cite{SURAWICZ2009AHA}, and acute ischaemia and infarction \cite{WAGNER2009AHA}. In this project, criteria of this kind are represented in an evidence schema described in Chapter 3. A diagnosis that depends on such criteria can only be verified when the corresponding leads are available.

Computational ECG analysis began with signal-processing algorithms and rule-based interpretation. The Pan--Tompkins algorithm, for example, detects QRS complexes through band-pass filtering, differentiation, squaring and adaptive thresholding \cite{PAN1985QRS}, and computer interpretation programs combine such measurements with diagnostic rules \cite{SCHLAPFER2017COMPUTER}. These methods are transparent but depend on hand-designed features and thresholds. Deep neural networks instead learn features directly from the waveform. Networks trained on large labelled datasets have achieved strong performance in rhythm and abnormality classification, including benchmark studies on PTB-XL \cite{HANNUN2019CARDIOLOGIST,STRODTHOFF2021PTBXL}. Because such models learn their own features, high accuracy does not by itself reveal which signal properties a model relies on.

Self-supervised learning reduces the dependence on labelled data by training an encoder on tasks constructed from unlabelled signals. ECG-FM combines wav2vec 2.0 signal masking, the contrastive multi-segment coding objective of CLOCS, and random lead masking, which removes leads during pretraining so that the encoder learns to operate on varying lead subsets \cite{MCKEEN2025ECGFM,OH2022LEADAGNOSTIC}. The resulting transformer-based model has about 90.9 million parameters and was pretrained on recordings from MIMIC-IV-ECG and the PhysioNet 2021 challenge collection \cite{MCKEEN2025ECGFM}.

Reduced-lead analysis has been studied both as a classification problem and as a reconstruction problem. The PhysioNet/Computing in Cardiology Challenge 2021 compared classification from several lead combinations \cite{REYNA2022ISSUES}, and lead-selection studies have shown that a small subset of leads can generalise well for some abnormality classes \cite{LAI2021LEADSELECTION}. Reconstruction approaches estimate the missing leads from those that remain \cite{PRESACAN2025RECON}. Both lines of work evaluate an end result, a classification score or a reconstructed waveform, rather than whether the evidence for each diagnostic claim was observable in the recorded leads.

It is therefore necessary to distinguish three properties of a model under lead removal. Classification performance, measured by metrics such as the area under the receiver operating characteristic curve, indicates agreement between predicted and reference labels. Representation similarity, such as the cosine similarity between the representations of a degraded and a full recording, indicates how much the encoder's output changes. Evidence-grounded reasoning concerns whether the claims made to justify a prediction can be verified from the leads that are present. The first two properties can remain high while the third fails, because neither of them examines the evidence behind individual claims. A review of explainable AI in cardiology found that many studies do not evaluate their explanations at all \cite{SALIH2024XAICARDIO}.

Assessing grounding requires two additional ingredients. The first is a set of measurable physiological findings, such as heart rate, rhythm regularity, interval durations and ST-segment deviations, computed only from the available leads. The second is structured clinical knowledge that states which findings and which leads are required for each diagnosis. SNOMED CT is a clinical terminology whose encoded concepts can be combined with health records in clinical decision-support systems \cite{SNOMED2026CDSGUIDE}. Comparable three-way judgements are established in automated fact verification for text: the FEVER dataset labels claims as Supported, Refuted or NotEnoughInfo with respect to evidence \cite{THORNE2018FEVER}. The UNAVAILABLE state used in this project is a proposed adaptation of this distinction to ECG claims, in which lead availability determines whether the required evidence can be assessed.

\section{PROBLEM DEFINITION}
The problem addressed in this project is to determine how the progressive removal of ECG leads changes the physiological evidence available to an ECG foundation model, and whether the diagnostic claims of a reasoning system built on that model remain grounded in the signals that survive. Three evidence states are distinguished for each claim. A claim is SUPPORTED when the leads required to assess it are present and the measured physiological evidence agrees with it. A claim is CONTRADICTED when the required leads are present but the measured evidence conflicts with it. A claim is UNAVAILABLE when the evidence needed to evaluate it cannot be assessed because one or more required leads are absent. The third state is kept separate so that missing information is never counted as either confirmation or refutation.

A model may retain similar internal representations, or produce a confident prediction, even when the evidence for a particular claim is no longer observable. Several factors may contribute. Some removed leads, such as lead III and the augmented leads, are linear combinations of leads I and II, so their removal hides rather than destroys information. A foundation model pretrained with random lead masking may also be designed to produce stable outputs from partial inputs \cite{MCKEEN2025ECGFM,OH2022LEADAGNOSTIC}. In addition, a classifier may rely on correlated features, such as heart rate or global morphology, that remain visible after the leads needed for a specific criterion have been removed. These explanations are hypotheses to be examined, not established findings.

To study when reasoning ceases to be grounded, the project defines the Reasoning Failure Boundary (RFB). The lead configurations are ordered from 12 leads to 1 lead, and for each configuration a grounding quality is computed as the proportion of supported claims among all claims that can be assessed, that is, the supported and contradicted claims. The RFB is the first configuration in this order at which grounding quality falls below a chosen threshold. The project's internal evaluation protocol, which is hash-verified but not externally registered (Section 3.3.12), uses a threshold of 0.60, with 0.50 and 0.70 as sensitivity values, and reports the transition from the last acceptable to the first failing configuration. Because the configurations are not nested subsets of one another, a recovery of grounding quality at a configuration with fewer leads is reported separately rather than forcing the boundary to be monotonic. The threshold is an analysis parameter selected by the project, not a clinically validated standard. The definition is fixed in the protocol, but the grounding gate and the procedure for extracting verifiable claims from the reasoning traces have not yet been implemented, so no RFB value has been measured at this stage.

Formally, the problem is stated as follows. Given 12-lead ECG recordings subjected to a fixed sequence of non-generative lead removals from 12 to 1 leads, the project investigates how the loss of observable physiological evidence relates to the representations of a frozen ECG foundation model, to the classification behaviour of a reasoning model built on those representations, and to the grounding of the clinical claims that the reasoning model makes, with the aim of identifying the conditions under which such claims remain supported by evidence that is actually present in the available leads.

\section{ORGANIZATION OF THE REPORT}
This report is organised into seven chapters. Chapter 1 presents the motivation, background, problem definition and organisation of the report. Chapter 2 reviews the literature on deep learning for ECG analysis, ECG foundation models, reduced-lead analysis, explainability and evidence verification, identifies the research gaps and states the research objectives. Chapter 3 describes the proposed methodology: the dataset pipeline, preprocessing, the seven lead-removal configurations, the frozen foundation model, the observation extractor and the Control Reasoning Transformer, together with the grounding gate, grounding metrics and Clinical Knowledge Prior, which are specified but not yet implemented. Chapter 4 describes the datasets, the hardware and software environment and the experiments conducted, and analyses the results obtained so far with their limitations. Chapter 5 discusses social impact and sustainability, Chapter 6 presents the conclusions and future work, and Chapter 7 lists the references.